\documentclass[11pt]{article}
\makeatletter\def\input@path{{../estilo/}}\makeatother
\usepackage{didistribuidores}
\renewcommand{\documentocorto}{Reporte final}

\begin{document}
\portada{Reporte final del proyecto}{Pedidos que no se pierden ni se cobran dos veces}%
  {Un sistema distribuido de pedidos y entrega de comida, construido de monolito a servicios en cuatro etapas}%
  {Ramas \cod{01-monolito}, \cod{02-separacion}, \cod{03-separacion} y \cod{04-resiliencia}}{28 de noviembre de 2026}

\tableofcontents
\newpage

\section{El problema}

Cuando alguien pide comida a domicilio en hora pico, las aplicaciones se ponen lentas, pierden pedidos
o cobran dos veces. Muchas veces la causa es la propia red: la aplicación no recibe la respuesta,
reintenta, y el servidor procesa el mismo pedido otra vez.

\textbf{Para quién.} El usuario principal es el \textbf{cliente}: quiere pedir, pagar una sola vez y
recibir su comida. Los secundarios son el \textbf{repartidor}, que acepta y entrega, y el
\textbf{restaurante}, que prepara (en este alcance se asume que prepara todo pedido pagado).

\textbf{Qué priorizamos.} La \textbf{consistencia sobre la disponibilidad}: preferimos retrasar o
rechazar un pedido antes que perderlo, duplicarlo o cobrarlo dos veces. Las garantías que nunca se rompen:

\begin{itemize}
  \item Un reintento con la misma \cod{Idempotency-Key} devuelve el mismo pedido, sin segundo cobro.
  \item Nunca se cobra un pedido que no se confirmó, y nunca se pierde uno ya confirmado.
  \item La entrega se registra una sola vez, aunque el repartidor la confirme dos veces.
  \item Si nadie toma un pedido pagado a tiempo, se cancela y se reembolsa, una sola vez.
  \item Los avisos al cliente llegan una vez cada uno, aunque el servicio de avisos haya estado apagado.
\end{itemize}

\textbf{Por qué es un sistema distribuido.} Un aviso al cliente o la búsqueda de repartidor no deben
impedir que un pedido se confirme; por eso \cod{notificaciones} y \cod{repartidores} corren en procesos
aparte, cada uno con su base. \cod{repartidores} no guarda estado en memoria, así que corre en dos copias.
Al comunicarse por la red aparecen los problemas del curso: tiempos límite, reintentos, duplicados,
desorden y caídas.

\section{Diagrama funcional}

Los tres flujos principales, de punta a punta: qué hace la persona y qué partes del sistema intervienen.

\begin{center}
% Diagrama funcional del reporte final (se incluye con \input)
\begin{tikzpicture}[x=1cm,y=-1cm, every node/.style={font=\sffamily\scriptsize},
  msg/.style={-Stealth, thick, azulsd}, res/.style={-Stealth, thick, verdeok, dashed},
  et/.style={fill=white, inner sep=1.5pt},
  nota/.style={fill=grisfondo, draw=grislinea, rounded corners=2pt, align=left, inner sep=3pt},
  flujo/.style={fill=azulsd!8, draw=azulsd!40, rounded corners=2pt, anchor=west, text width=15.2cm, inner sep=3pt}]
  \foreach \x/\n in {0/Cliente, 2.7/Gateway, 5.4/Núcleo, 8.6/Repartidores, 11.0/RabbitMQ, 13.4/Notificaciones} {
    \node[font=\sffamily\scriptsize\bfseries, text=azulsd] at (\x,0) {\n};
    \draw[grislinea, dashed] (\x,0.35) -- (\x,15.5);
  }
  % Flujo 1
  \node[flujo] at (-0.9,0.9) {\textbf{1. Pedir y pagar}};
  \draw[msg] (0,1.7) -- node[et,above]{POST /api/pedidos} (2.7,1.7);
  \draw[msg] (2.7,1.9) -- (5.4,1.9);
  \node[nota, anchor=north west, text width=2.8cm] at (5.55,2.05) {valida el menú, calcula el total; \texttt{UNIQUE} sobre la clave};
  \draw[res] (5.4,3.5) -- node[et,above]{201 (repetido: 200, el mismo)} (0,3.5);
  \draw[msg] (0,4.2) -- node[et,above]{POST .../pago} (2.7,4.2);
  \draw[msg] (2.7,4.4) -- (5.4,4.4);
  \node[nota, anchor=north west, text width=2.8cm] at (5.55,4.55) {una transacción: PAGADO, cobro y evento en \texttt{outbox}};
  \draw[msg] (5.4,6.1) -- node[et,above,pos=0.62]{publicador, con confirmación} (11.0,6.1);
  \draw[msg] (11.0,6.5) -- node[et,above]{entrega} (13.4,6.5);
  \node[nota, anchor=north east, text width=2.4cm] at (13.25,6.7) {guarda el aviso y luego \emph{ack}};
  % Flujo 2
  \node[flujo] at (-0.9,7.9) {\textbf{2. Aceptar y entregar} (el repartidor entra por la misma puerta)};
  \draw[msg] (0,8.7) -- node[et,above]{POST .../aceptar} (2.7,8.7);
  \draw[msg] (2.7,8.9) -- (5.4,8.9);
  \draw[msg] (5.4,9.4) -- node[et,above]{¿quién es? (1 s / 3 s, breaker)} (8.6,9.4);
  \draw[res] (8.6,9.9) -- node[et,below]{zona y disponibilidad} (5.4,9.9);
  \node[nota, anchor=north west, text width=2.8cm] at (5.55,10.45) {\texttt{UPDATE} condicional: gana el primero; el otro recibe 409};
  \draw[msg] (0,12.0) -- node[et,above]{POST .../entrega + código} (2.7,12.0);
  \draw[msg] (2.7,12.2) -- (5.4,12.2);
  \node[nota, anchor=north west, text width=2.8cm] at (5.55,12.35) {ENTREGADO una sola vez (versión 4) y su evento};
  % Flujo 3
  \node[flujo] at (-0.9,13.9) {\textbf{3. Ver los avisos}};
  \draw[msg] (0,14.6) -- node[et,above]{GET /api/notificaciones/pedido/\{id\}} (2.7,14.6);
  \draw[msg] (2.7,14.8) -- (13.4,14.8);
  \draw[res] (13.4,15.2) -- node[et,below]{avisos ordenados por versión} (0,15.2);
\end{tikzpicture}

\end{center}

\begin{nota}[Lo que pasa sin que nadie haga nada]
Una tarea del núcleo revisa cada 5 segundos los pedidos: a los 30 s de pagado sin repartidor avisa al
cliente del retraso; a los 60 s lo cancela y reembolsa en una sola transacción; un pedido sin pagar se
cancela a los 30 s (valores de demostración; en producción, 10 y 20 minutos).
\end{nota}

\section{Diagrama de arquitectura}

\begin{center}
% Diagrama de arquitectura del reporte final (se incluye con \input)
\begin{tikzpicture}[every node/.style={font=\sffamily\small},
  svc/.style={draw=azulsd, fill=grisfondo, rounded corners=4pt, minimum height=1.25cm, align=center, text width=3.1cm},
  base/.style={draw=grisTexto, fill=white, shape=cylinder, shape border rotate=90, aspect=0.25, minimum height=1.0cm, minimum width=2.9cm, align=center, font=\sffamily\scriptsize},
  fl/.style={-Stealth, thick, azulsd}, et/.style={font=\sffamily\scriptsize, fill=white, inner sep=1.5pt}]
  \node[svc, text width=12.6cm, fill=white, draw=grisTexto] (cli) at (0,0) {Clientes: \texttt{interfaz.py}, panel web \texttt{/panel/}, Swagger, \texttt{curl}};
  \node[svc, text width=6.2cm, fill=azulsd!12] (gw) at (0,-1.9) {\textbf{Traefik} (gateway), \texttt{:8080}: única entrada\\ \scriptsize enruta por ruta y reparte entre copias};
  \node[svc] (nuc) at (-4.7,-4.1) {\textbf{Núcleo}\\ \scriptsize usuarios, restaurantes, pedidos, pagos, \texttt{outbox}, breaker};
  \node[svc] (rep) at (0,-4.1) {\textbf{Repartidores}\\ \scriptsize 2 copias, sin estado};
  \node[svc] (not) at (4.7,-4.1) {\textbf{Notificaciones}\\ \scriptsize consumidor y API};
  \node[svc, fill=ambar!10, draw=ambar] (rmq) at (2.35,-6.3) {\textbf{RabbitMQ}\\ \scriptsize cola durable y DLQ};
  \node[base] (b1) at (-4.7,-8.5) {didistribuidores};
  \node[base] (b2) at (0,-8.5) {repartidores\_db};
  \node[base] (b3) at (4.7,-8.5) {notificaciones\_db};
  \node[draw=grisTexto, dashed, rounded corners, fit=(b1)(b2)(b3), inner sep=7pt,
        label={[font=\sffamily\scriptsize,text=grisTexto]below:PostgreSQL 17: una instancia, una base lógica por servicio}] {};
  \draw[fl] (cli) -- node[et,right]{HTTP} (gw);
  \draw[fl] (gw.south west) -- (nuc.north); \draw[fl] (gw) -- (rep); \draw[fl] (gw.south east) -- (not.north);
  \draw[fl, dashed] (nuc) -- (rep);
  \node[et, align=center] at (-2.35,-5.05) {HTTP directo\\ (DNS de Docker)};
  \draw[fl] ([xshift=6mm]nuc.south) |- node[et,pos=0.75,below]{eventos del \texttt{outbox}} (rmq.west);
  \draw[fl] (rmq.east) -| (not.south);
  \draw[fl] ([xshift=-6mm]nuc.south) -- ([xshift=-6mm]b1.north);
  \draw[fl] (rep) -- (b2);
  \draw[fl] ([xshift=8mm]not.south) -- ([xshift=8mm]b3.north);
\end{tikzpicture}

\end{center}

\subsection{Dónde vive cada pieza}
\begin{tabularx}{\linewidth}{@{}P{3cm} Y Y@{}}
\toprule
\textsf{\textbf{Pieza}} & \textsf{\textbf{Dónde vive}} & \textsf{\textbf{Si falla}} \\
\midrule
Gateway & Traefik v3.6, configurado por etiquetas de Docker. Rutas: \cod{/api/usuarios}, \cod{/api/restaurantes}, \cod{/api/pedidos}, \cod{/api/pagos}, \cod{/salud}, \cod{/estado}, \cod{/panel} al núcleo; \cod{/api/repartidores} y \cod{/api/notificaciones} a sus servicios & Nadie entra desde afuera, pero los servicios siguen hablándose y no se pierde nada (bitácora 4, prueba 7). Punto único de falla aceptado \\
Directorio & DNS de Docker Compose: los servicios se buscan por nombre (\cod{repartidores}, \cod{rabbitmq}, \cod{db}) & El núcleo resuelve el nombre con tiempo límite de 1 s; si no se resuelve, 503 y el breaker se abre \\
Balanceo & Traefik reparte por turnos entre las 2 copias de repartidores y saca a la que no contesta \cod{/salud}; las llamadas internas se reparten por el DNS & Con una copia caída, la otra atiende todo; con las dos, entra el breaker \\
Datos & Una instancia de PostgreSQL con tres bases lógicas; cada servicio solo toca la suya & Durante un reinicio hay errores; después los pools se reconectan solos. Si cae para siempre, todo se detiene (declarado) \\
Operaciones entre servicios & Eventos: \cod{outbox} + RabbitMQ (al menos una vez + consumidor idempotente). Consulta a repartidores: HTTP con tiempo límite, un reintento seguro y circuit breaker & Ver la sección~\ref{sec:fallas} \\
\bottomrule
\end{tabularx}

\section{Cómo se construyó: cuatro etapas}

\begin{tabularx}{\linewidth}{@{}l P{3.3cm} Y l@{}}
\toprule
\textsf{\textbf{Rama}} & \textsf{\textbf{Qué se agregó}} & \textsf{\textbf{Qué demuestra}} & \textsf{\textbf{Bitácora}} \\
\midrule
\cod{01-monolito} & Un proceso, una base, módulos con tres capas & Idempotencia con \cod{UNIQUE}, cambios de estado con \cod{UPDATE} condicional, pago y reembolso en una transacción, relojes, \cod{outbox} & 1 \\
\cod{02-separacion} & \cod{notificaciones} aparte, RabbitMQ & Los pedidos funcionan con los avisos apagados; ningún aviso se pierde ni se duplica; DLQ; orden por versión & 2 \\
\cod{03-separacion} & \cod{repartidores} aparte (2 copias), Traefik & Una sola entrada, balanceo, base propia, tiempo límite (503), aviso de retraso & 3 \\
\cod{04-resiliencia} & Breaker, reintentos, huella, \cod{/estado}, interfaces & Fallar rápido, recuperarse solo, rechazar claves reutilizadas; matriz de caídas con Docker real & 4 \\
\bottomrule
\end{tabularx}

Separar por etapas funcionó porque desde la etapa 1 cada módulo solo hablaba con otro por su
\cod{servicio.py} y cada cambio de estado ya dejaba su evento: separar un servicio fue cambiar una llamada a
función por una llamada de red, sin reescribir la lógica de pedidos.

\section{Decisiones}

{\small\begin{longtable}{@{}P{3.0cm} P{3.4cm} P{3.3cm} P{5.2cm}@{}}
\toprule
\textsf{\textbf{Decisión}} & \textsf{\textbf{Qué dejamos atrás}} & \textsf{\textbf{Alternativa descartada}} & \textsf{\textbf{Por qué}} \\
\midrule
\endhead
Consistencia sobre disponibilidad & Aceptar cualquier pedido aunque algo falle; velocidad en hora pico & Responder siempre ``sí'' y corregir después & Preferimos un 503 claro a un pedido en el limbo o un doble cobro \\
Idempotencia con \cod{UNIQUE} en la base & Simplicidad de revisar en código & Un diccionario en memoria que revisa si ya existe & Entre revisar y guardar hay un hueco; la base lo cierra de forma atómica (5 peticiones simultáneas: 1 pedido) \\
\cod{pagos} dentro del núcleo & Un servicio de pagos independiente & Pagos como servicio con saga y compensación & Cobrar, cambiar el estado y reembolsar son una transacción local: no puede haber cobro sin pedido \\
\cod{outbox} + RabbitMQ para los avisos & La simplicidad de una llamada directa & Llamar a notificaciones por HTTP dentro de la petición & Si los avisos fallan, el pedido no se bloquea ni se deshace; el evento nace en la misma transacción \\
\emph{ack} después de guardar & Un poco de rendimiento & \emph{ack} automático al recibir & Con \emph{ack} temprano, una caída entre recibir y guardar perdería el aviso; así lo peor es un duplicado, que se ignora \\
Repartidores informa, pedidos decide & Que repartidores asigne pedidos & Que el servicio de repartidores elija quién gana & La regla de quién gana vive con el dueño del estado: un \cod{UPDATE} condicional, sin cruzar la red \\
503 si repartidores no responde & Disponibilidad al aceptar & Aceptar a ciegas & Sin preguntar no se sabe si el repartidor existe o es de esa zona \\
Circuit breaker hecho a mano, solo hacia repartidores & Una librería probada & Una librería de resiliencia & Es la única llamada síncrona entre servicios; a mano se explica y se prueba línea por línea \\
Una instancia de PostgreSQL con tres bases & Tolerancia a fallas de la base & Una instancia por servicio, o NoSQL & Cada servicio es dueño de sus datos sin multiplicar contenedores; el punto único de falla queda declarado \\
Traefik y DNS de Compose & El ecosistema del lab (Java) & Spring Cloud Gateway y Eureka & Usamos Python; Docker ya resuelve nombres y Traefik se configura con etiquetas \\
\bottomrule
\end{longtable}}

\section{Manejo de fallas}\label{sec:fallas}

Todas las filas se comprobaron con el sistema real; la columna de evidencia apunta a la bitácora y la
prueba. Las de la etapa 4 se corrieron en Docker Desktop el 18 de noviembre de 2026.

{\small\begin{longtable}{@{}P{3.0cm} P{6.0cm} P{3.9cm} P{1.9cm}@{}}
\toprule
\textsf{\textbf{Qué se cae}} & \textsf{\textbf{Qué pasa}} & \textsf{\textbf{Qué lo resuelve}} & \textsf{\textbf{Evidencia}} \\
\midrule
\endhead
Notificaciones & Los pedidos se crean, pagan y entregan; los avisos esperan en la cola y llegan al volver, una vez cada uno & Cola durable, \emph{ack} tardío, \cod{evento\_id} único & B2, p. 2 \\
RabbitMQ & Los pedidos siguen; los eventos esperan en el \cod{outbox} y se publican al volver & Outbox y publicador que reintenta & B2, p. 6 \\
RabbitMQ se reinicia & Los mensajes en cola sobreviven & Cola durable, mensajes persistentes, \cod{hostname} fijo & B2, p. 7 \\
Una copia de repartidores & Nada se nota desde afuera & Traefik la saca del reparto & B3, p. 5 \\
Las dos copias & Crear y pagar siguen; aceptar da 503 en 1 s y, tras 3 fallos, al instante; aviso de retraso; cancelación y reembolso al vencer & Tiempo límite, breaker, relojes & B3 p. 6; B4 p. 4 \\
Repartidores lento & Tres intentos de 3 s y luego 503 en 7 ms; se recupera solo & Tiempo límite y breaker & B4, p. 2 \\
Una copia congelada & Algunas llamadas internas tardan 3 s; el breaker puede abrirse & Parcial (limitación declarada) & B4, p. 3 \\
El núcleo muere tras un pago & El evento sale al volver, una sola vez & Outbox en la misma transacción & B4, p. 5 \\
PostgreSQL se reinicia & Errores unos segundos; luego todo se recupera sin reiniciar servicios & Pools con \cod{check\_connection} & B4, p. 6 \\
El gateway & Nadie entra; adentro siguen hablándose & Llamada interna directa & B4, p. 7 \\
Mensaje inválido o duplicado & Inválido a la DLQ sin bloquear; duplicado ignorado & \emph{nack} sin reencolar; \cod{UNIQUE} & B2, p. 3--4 \\
Petición repetida o simultánea & Mismo pedido; un solo cobro; un solo ganador & \cod{UNIQUE} y \cod{UPDATE} condicional & B1, secc. 5--8 \\
Clave reutilizada con otro pedido & 422 & Huella SHA-256 & B4, p. 1 \\
\bottomrule
\end{longtable}}

\begin{nota}[Una falla que encontramos y corregimos con las pruebas reales]
Con repartidores detenido en Docker Desktop para Windows, cada intento tardaba 8 s: resolver un nombre
inexistente toma unos 4 s y el tiempo límite de \cod{httpx} no cubre esa parte (el reintento la
duplicaba). Ahora el nombre se resuelve con su propio límite de 1 s y sin reintento: de 8 s a 1 s, y
7 ms con el circuito abierto (bitácora 4, sección de hallazgos).
\end{nota}

\section{Cómo levantarlo y probarlo}

\begin{comando}
git clone https://github.com/YukiGab/sd-2027-1-didistribuidores
cd sd-2027-1-didistribuidores/Proyecto
docker compose up --build -d
docker compose ps                       # siete contenedores
curl http://localhost:8080/salud        # {"estado":"ok"}
curl http://localhost:8080/estado       # circuito CERRADO, 0 eventos pendientes
python interfaz/interfaz.py             # aplicación de escritorio (solo biblioteca estándar)
\end{comando}

\begin{itemize}
  \item Panel web: \url{http://localhost:8080/panel/} · Swagger: \url{http://localhost:8080/docs} ·
        Traefik: \url{http://localhost:8081} · RabbitMQ: \url{http://localhost:15672} (\cod{sd}/\cod{sd}).
  \item Para empezar de cero: \cod{docker compose down -v} (los scripts de \cod{db/} solo corren con el
        volumen vacío).
  \item Las pruebas de cada etapa, con comandos y resultados: \cod{documentacion/pruebas/} (PDF y Word).
        Las de la etapa 4 y la regresión son scripts que se pueden repetir tal cual.
  \item El uso paso a paso, con capturas: \cod{documentacion/manual-de-usuario/}.
\end{itemize}

\textbf{Demostración sugerida (5 minutos).} 1) Pedir y pagar como cliente en la interfaz y repetir la
petición: mismo pedido. 2) Aceptar y entregar como repartidor con el código; ver los tres avisos.
3) \cod{docker compose stop notificaciones}, hacer un pedido, encenderlo y ver llegar los avisos.
4) \cod{docker compose stop repartidores}: el circuito se pone en rojo y las respuestas son instantáneas;
encenderlo y verlo cerrarse solo.

\section{Qué falta y qué haríamos distinto}

\begin{tabularx}{\linewidth}{@{}P{4.2cm} Y@{}}
\toprule
\textsf{\textbf{Qué falta}} & \textsf{\textbf{Cómo se resolvería}} \\
\midrule
PostgreSQL y Traefik son puntos únicos de falla & Réplicas de la base con failover y varias copias del gateway \\
Las llamadas internas no usan los \emph{healthchecks}: una copia congelada recibe tráfico interno & Llamar a repartidores a través del gateway o con un directorio que conozca la salud \\
El breaker vive en cada proceso del núcleo & Un breaker compartido (por ejemplo, en Redis) si hubiera varias copias del núcleo \\
Sin autenticación: el usuario se identifica con su id y \cod{GET /api/pedidos/\{id\}} muestra el código de entrega a quien conozca el id & Autenticación con tokens; mostrar el código solo al dueño del pedido. (Sí se quitó el código de la lista que ve el repartidor.) \\
Los mensajes de la cola de fallidos no se reprocesan & Una herramienta para revisarlos y reenviarlos \\
Los avisos solo se guardan & Envío real por correo o notificación push \\
La base del núcleo se llama \cod{didistribuidores} y el PRD dice \cod{pedidos\_db} & Renombrarla en una versión siguiente \\
\bottomrule
\end{tabularx}

\textbf{Qué haríamos distinto.} Correr las pruebas contra Docker real desde el primer commit de cada etapa:
el hallazgo del DNS solo apareció en Windows y al final. Automatizar esas pruebas en un pipeline (como los
scripts de la etapa 4) en lugar de capturas manuales. Y definir desde el inicio una interfaz mínima, que
ayudó mucho a explicar el sistema pero llegó en la última etapa.

\section{Quién hizo qué}

El reparto se acordó desde el inicio: \textbf{Erick} (Persona A) con las etapas 1 y 3 y \textbf{Yuki}
(Persona B) con las etapas 2 y 4. Cada pull request lo revisó la otra persona antes de mezclarlo.

\begin{tabularx}{\linewidth}{@{}P{4.2cm} Y@{}}
\toprule
\textsf{\textbf{Integrante}} & \textsf{\textbf{Aportes}} \\
\midrule
Erick Isaac Echeverria Goicochea (@ErickI07) & Etapa 1: módulos, esquema, idempotencia, pagos, aceptación, entrega, relojes con reembolso y sus pruebas. Etapa 3: servicio de repartidores con base propia y dos copias, Traefik, cliente HTTP con tiempo límite, aviso de retraso y sus 10 pruebas. Panel web de demostración. Bitácoras originales de las etapas 1 a 3. Revisión de las etapas 2 y 4. \\
Yukioayax Canek Gabriel Hernández (@YukiGab) & Etapa 1: Docker Compose, datos semilla y tabla \cod{outbox}. Etapa 2: servicio de notificaciones, publicador del \cod{outbox}, RabbitMQ con DLQ y sus 7 pruebas. Etapa 4: circuit breaker, reintentos seguros, huella de la petición, \cod{/estado}, recuperación de los pools, pruebas reales de caídas, corrección del límite de la resolución de nombres e \cod{interfaz.py}. Documentación final: bitácoras en \LaTeX, pruebas, manual y este reporte. Revisión de las etapas 1 y 3. \\
Los dos & PRD y decisiones de alcance (acordadas con el profesor), contrato de eventos, revisión cruzada de todos los pull requests y preparación de la defensa. \\
\bottomrule
\end{tabularx}

\end{document}
